
\section{Results and Discussion}\label{sec:results}

The strongest result is conceptual rather than promotional: the option-only Markowitz
object is coherent, auditable, and measurable on the held-out panel. The empirical
investment result is more conditional. Gross option portfolios can look attractive, but
gross Sharpe is a research diagnostic rather than a trading result. VIX options can improve
the design as volatility instruments and hedge candidates only when exact VRO/SOQ settlement
coverage supports the reported rows.

The body keeps the central evidence in view: out-of-sample performance in
Table~\ref{tab:performance}, post-cost survival in Table~\ref{tab:post_cost_survival},
factor exposure in Table~\ref{tab:factor_regression}, the realized path in
Figure~\ref{fig:growth}, the random feasible comparison in Figure~\ref{fig:random_sharpe},
the VIX settlement gate in Table~\ref{tab:vix_settlement}, claim strength in
Table~\ref{tab:claim_strength}, and the claim audit in Table~\ref{tab:claim_audit}.
Detailed exposure, liquidity, attribution, regime, forecast-ablation, simulation, and
multiple-testing diagnostics appear in Appendix A.

\subsection{Headline evidence}

\begin{table}[H]
\centering
\scriptsize
\IfFileExists{tables/portfolio_performance.tex}
  {\resizebox{\textwidth}{!}{\begin{tabular}{llrrrrrrr}
\toprule
Return basis & Strategy & Ann. return & Ann. vol & Sharpe & Sortino & Calmar & Omega & Info. ratio \\
\midrule
Gross before costs & Equity-option Greek Markowitz & 0.894 & 1.062 & 0.842 & 1.770 & 1.020 & 1.888 & 0.598 \\
Gross before costs & Greek Markowitz + VIX & 1.246 & 0.907 & 1.374 & 3.385 & 2.034 & 2.784 & 1.200 \\
Gross before costs & Beta/delta-neutral + VIX & 1.172 & 0.829 & 1.414 & 3.410 & 2.057 & 2.948 & 1.147 \\
Gross before costs & Cost-aware Sortino + VIX & 5.020 & 4.949 & 1.014 & 3.349 & 4.993 & 2.471 & 0.971 \\
Gross before costs & Equal premium & -0.056 & 1.878 & -0.030 & -0.055 & -0.056 & 0.977 & -0.120 \\
Gross before costs & Equal risk & 0.104 & 1.275 & 0.082 & 0.146 & 0.107 & 1.064 & -0.058 \\
Gross before costs & VIX hedge sleeve & -9.285 & 1.444 & -6.429 & -3.089 &  & 0.040 & -6.132 \\
Gross before costs & Delta-matched equities & 1.095 & 0.931 & 1.176 & 2.317 & 1.225 & 2.367 & 1.041 \\
Gross before costs & Underlying Markowitz & 0.179 & 0.196 & 0.909 & 1.856 & 0.590 & 1.929 & 0.072 \\
Post-cost research & Equity-option Greek Markowitz & -7.900 & 5.476 & -1.443 & -1.347 & -6.357 & 0.098 & -1.545 \\
Post-cost research & Greek Markowitz + VIX & -6.079 & 4.827 & -1.259 & -1.198 & -5.136 & 0.116 & -1.345 \\
Post-cost research & Beta/delta-neutral + VIX & -3.896 & 2.841 & -1.371 & -1.300 & -3.895 & 0.151 & -1.508 \\
Post-cost research & Cost-aware Sortino + VIX & -157.425 & 148.976 & -1.057 & -1.019 & -1.985 & 0.016 & -1.003 \\
Post-cost research & Equal premium & -3.165 & 2.378 & -1.331 & -1.502 & -3.051 & 0.359 & -1.372 \\
Post-cost research & Equal risk & -1.276 & 1.367 & -0.934 & -1.199 & -1.276 & 0.505 & -1.082 \\
Post-cost research & VIX hedge sleeve & -14.364 & 4.209 & -3.413 & -2.445 & -12.391 & 0.016 & -3.357 \\
Post-cost research & Delta-matched equities & 1.095 & 0.931 & 1.176 & 2.317 & 1.225 & 2.367 & 1.041 \\
Post-cost research & Underlying Markowitz & 0.179 & 0.196 & 0.909 & 1.856 & 0.590 & 1.929 & 0.072 \\
\bottomrule
\end{tabular}
}}
  {\begin{tabular}{llll}\toprule Strategy & Return & Sharpe & Claim status \\\midrule
  Missing & Run the empirical pipeline & -- & -- \\\bottomrule\end{tabular}}
\caption{Gross and stress-cost out-of-sample performance. Gross rows test the option-only
Markowitz object before frictions. The ``Post-cost research'' rows apply the
\emph{conservative stress-cost layer}: full spread crossing on every leg, explicit fees,
slippage, borrow, capacity penalties, simulated margin funding, and assignment-risk
penalties, compounded monthly. This layer is a stress bound, not an execution estimate;
the milder \emph{executable-cost scenarios} appear in
Table~\ref{tab:post_cost_survival}. Annualized returns are arithmetic monthly means
scaled by twelve, so stressed rows can fall below \(-100\%\) per year
(Section~\ref{sec:metrics} records the conventions).}
\label{tab:performance}
\end{table}

Table~\ref{tab:performance} separates two questions. The gross rows ask whether the
mathematical allocation object works on held-out option returns. The stress-cost rows ask
whether the same object would still look attractive if every implementation friction were
charged at its most conservative research setting simultaneously. Between those poles,
Table~\ref{tab:post_cost_survival} asks whether the result survives midpoint, half-spread,
and full-spread executable-cost scenarios with required-capital denominators, no-fill
diagnostics, and margin/stress-capital screens. Gross Sharpe is not treated as a trading
result under any of these readings.

On the gross ledger, the three optimizer-based option sleeves earn annualized returns of
\(0.894\), \(1.246\), and \(1.172\) (per unit of NAV) at annualized volatilities near
\(0.8\)--\(1.1\): out-of-sample Sharpe ratios of \(0.842\) for the equity-option Greek
Markowitz sleeve, \(1.374\) with the VIX overlay, and \(1.414\) for the beta/delta-neutral
variant. The circular-block-bootstrap 90\% confidence intervals are wide---\([0.083,
1.679]\), \([0.602, 2.222]\), and \([0.677, 2.220]\) respectively
(Table~\ref{tab:performance_diag_app} and the full ledger in
Table~\ref{tab:inference_app})---so the point estimates should be read as one draw from a
sixty-month window, not as stable population quantities. The naive allocations confirm
that construction, not option exposure per se, drives the result: equal-premium and
equal-risk books earn gross Sharpe ratios of \(-0.030\) and \(0.082\), and the always-long
VIX hedge sleeve loses its premium almost every month (Sharpe \(-6.429\), maximum drawdown
undefined because the compounded path absorbs at zero wealth).

The benchmark comparison is closer than the option-sleeve numbers alone suggest. The
delta-matched equity benchmark, which holds the same underlying names at the sleeves'
decision-date delta exposure without any options, earns a gross Sharpe of \(1.176\); the
plain underlying Markowitz portfolio earns \(0.909\). The option sleeves' advantage over
delta-matched equities is therefore roughly \(0.2\) Sharpe units at the point estimate,
well inside the bootstrap intervals. The framework's contribution is the accounting and
the risk decomposition, not a claim that these particular sleeves dominate their equity
shadows.

\subsection{Two cost layers}

The paper reports two deliberately different post-cost accountings, and they must not be
conflated. The \emph{executable-cost scenarios}
(Tables~\ref{tab:post_cost_survival} and~\ref{tab:capacity_market_impact_app}) charge
entry frictions only---quoted spread at midpoint, half, or full crossing plus explicit
fees---at an average monthly cost, under full-spread crossing, of roughly \(0.7\%\) to
\(3.2\%\) of NAV for the option sleeves (about \(8\%\) to \(38\%\) annualized). These
scenarios fail closed: a position rejected for missing quotes, capacity, or
assignment risk forfeits its gross return contribution instead of keeping it, so a
rejection can move net performance in either direction. The
\emph{conservative stress-cost layer} (the ``Post-cost research'' rows of
Table~\ref{tab:performance}, with components in Table~\ref{tab:costs_app}) additionally
charges round-trip spread on every rebalance leg, slippage, borrow, capacity penalties,
simulated margin funding on gross margin of roughly \(1\) to \(4.6\) times NAV, and
assignment-risk penalties. Its annualized cost drag is \(5.1\) to \(8.8\) NAV units per
year for the optimizer sleeves---one to two orders of magnitude larger than the executable
scenarios---which is why the same strategies show net Sharpe ratios of \(+0.91\),
\(+0.29\), and \(+0.68\) in Table~\ref{tab:post_cost_survival} but \(-1.44\), \(-1.26\),
and \(-1.37\) (with annualized returns as extreme as \(-790\%\) under the
arithmetic-mean convention) in Table~\ref{tab:performance}.

The reconciliation is a statement about claims, not about arithmetic. The executable-cost
scenarios are the realistic post-cost evidence for a small, patient allocator paying
quoted spreads once per month; under them the leading sleeves survive with full-spread
Sharpe ratios of \(0.913\), \(0.294\), and \(0.678\), while the naive books and the hedge
sleeve fail. Two features of these numbers deserve comment. The equity-option sleeve's
net Sharpe exceeds its gross because the trades the fill screens reject were, on net,
losers for that sleeve; and the VIX overlay pays the largest toll because its wing
positions cross the widest spreads and hit capacity screens most often. The stress-cost
layer is a bound: it shows that if margin funding, capacity,
and round-trip frictions are all charged at conservative research settings, no option
sleeve in this study survives. The realistic post-cost outcome likely lies between the two
layers and is scenario-dependent; the claim audit (Table~\ref{tab:claim_audit}) therefore records
post-cost viability as a conditional research finding rather than a robust result, and no
wording in this paper should be read as post-cost survival under all cost assumptions.

\subsection{Inference and multiple testing}

The inference layer keeps the gross numbers honest. In per-month units, the probabilistic
Sharpe ratios of the leading gross sleeves are high---\(0.979\), \(0.999\), and \(1.000\)
for the equity-option, VIX-overlay, and beta/delta-neutral variants---so each sleeve's
Sharpe is individually distinguishable from zero at conventional levels
(Table~\ref{tab:reality_check_app}). Two multiple-testing corrections then remove most of
that comfort. First, the centered reality-check family \(p\)-value across all evaluated
variants is \(0.078\): the observed maximum studentized statistic of \(3.161\) sits below
the bootstrap 95th percentile of \(3.424\), so the family-wise test fails the 5\% hurdle.
Second, the deflated Sharpe ratio is \(0.000\) for every gross variant and at most
\(0.51\) across the executable-cost variants. After
deflating by the dispersion of the evaluated variant set---which includes the
\(-6.4\)-Sharpe hedge sleeve and therefore a very wide cross-sectional spread---no
variant's Sharpe ratio is distinguishable from the expected maximum under the null of no
skill. These corrections support, rather than undermine, the paper's restrained thesis:
sixty months of data cannot certify skill for any variant in this family, and the paper
does not claim otherwise.

\begin{table}[H]
\centering
\scriptsize
\IfFileExists{tables/post_cost_survival.tex}
  {\resizebox{\textwidth}{!}{\begin{tabular}{lllllll}
\toprule
Strategy & Gross Sharpe & Net Sharpe & Net Calmar & Avg spread cost & Capacity used & Survives? \\
\midrule
Equity-option Greek Markowitz & 0.8421194565895301 & 0.7668067218645858 & 0.9148415644409728 & 0.004454581057950615 & 0.9901002345178692 & yes \\
Greek Markowitz + VIX & 1.3743885779509968 & 1.0258219183556512 & 1.331603783957096 & 0.02189086589088895 & 0.9988451122012506 & yes \\
Beta/delta-neutral + VIX & 1.4135236431740843 & 0.9508576292459209 & 1.2239969998479348 & 0.026468684884072495 & 0.9997049345720702 & yes \\
Equal premium & -0.02955676336526383 & -0.20462197552298175 & -0.38425395579926913 & 0.020340974748224574 & 0.9993803841618195 & no \\
Equal risk & 0.0818308658646486 & -0.11815498556113241 & -0.15161444631136373 & 0.015983091604746466 & 0.9990134907215874 & no \\
VIX hedge sleeve & -6.429184989935395 & -6.607601788466617 & nan & 0.022500000000000003 & 0.9885078059495911 & no \\
\bottomrule
\end{tabular}
}}
  {\begin{tabular}{llllll}\toprule Strategy & Gross Sharpe & Net Sharpe & Net Calmar & Cost & Survives \\\midrule
  Missing & Run the empirical pipeline & -- & -- & -- & -- \\\bottomrule\end{tabular}}
\caption{Post-cost survival under the full-spread executable-cost scenario. This layer
charges entry frictions only (quoted spread crossed in full plus fees); it is the milder
of the two cost layers and is distinct from the conservative stress-cost layer of
Table~\ref{tab:performance}. A strategy survives only if net performance remains positive
after these frictions together with capacity screens, settlement gates, and risk-capital
checks.}
\label{tab:post_cost_survival}
\end{table}

\subsection{Factor exposure and benchmarks}

\begin{table}[H]
\centering
\scriptsize
\IfFileExists{tables/factor_regression_headline.tex}
  {\resizebox{\textwidth}{!}{\begin{tabular}{lrrrrrrrr}
\toprule
Strategy & Ann. alpha & Alpha HAC t & $R^2$ & Residual ann. vol & Beta SPY & Beta VX front & Beta dVIX & Beta dVVIX \\
\midrule
Equity-option Greek Markowitz & 0.405 & 1.005 & 0.618 & 0.634 & 0.005 & -1.368 & 0.046 & 0.003 \\
Greek Markowitz + VIX & 0.675 & 1.814 & 0.513 & 0.616 & 1.908 & -0.767 & 0.056 & -0.001 \\
Beta/delta-neutral + VIX & 0.845 & 2.425 & 0.429 & 0.601 & 0.715 & -0.522 & 0.038 & -0.001 \\
Cost-aware Sortino + VIX & 4.091 & 2.422 & 0.451 & 3.650 & 13.774 & -4.302 & 0.195 & 0.008 \\
Equal premium & -0.416 & -0.499 & 0.253 & 1.647 & 2.522 & -0.749 & 0.032 & -0.001 \\
Equal risk & -0.183 & -0.294 & 0.253 & 1.115 & 2.041 & -0.379 & 0.020 & -0.003 \\
VIX hedge sleeve & -8.807 & -12.880 & 0.252 & 1.264 & 3.562 & -1.180 & 0.093 & -0.013 \\
Delta-matched equities & 0.000 & 1.990 & 1.000 & 0.000 & -0.000 & -0.000 & -0.000 & 0.000 \\
Underlying Markowitz & 0.000 & 2.378 & 1.000 & 0.000 & -0.000 & -0.000 & -0.000 & 0.000 \\
\bottomrule
\end{tabular}
}}
  {\begin{tabular}{llll}\toprule Strategy & Alpha & Beta controls & HAC t-stat \\\midrule
  Missing & Run the empirical pipeline & -- & -- \\\bottomrule\end{tabular}}
\caption{Headline factor regressions with HAC (Newey--West) standard errors. The table
asks whether option returns are explained by equity drift or volatility state controls
before any alpha language is considered; the full eight-underlying regression appears in
Appendix A (Table~\ref{tab:factor_regression_full_app}). The two equity benchmark rows are
exact linear combinations of the regressors (\(R^2=1.000\), zero residual variance), so
their alpha \(t\)-statistics are numerical artifacts and should be ignored.}
\label{tab:factor_regression}
\end{table}

The regression is an audit, not a causal proof. Against SPY, the eight underlying
equities, VX front returns, and VIX/VVIX changes, the annualized alphas of the three
optimizer sleeves are \(0.405\) (HAC \(t=1.005\)), \(0.675\) (\(t=1.814\)), and \(0.845\)
(\(t=2.425\)), with \(R^2\) of \(0.618\), \(0.513\), and \(0.429\) and residual annualized
volatility near \(0.6\)---the factor block explains roughly half of the option-book
variance and leaves an economically large unexplained remainder
(Table~\ref{tab:factor_regression_full_app}). The \(t\)-statistics carry the
small-sample \(n/(n-k)\) HAC correction and are referred to a \(t(n-k)\) distribution.
Only the beta/delta-neutral variant clears
a conventional two-sided 5\% threshold, and none of these point estimates survives the
family-wise corrections above, so the paper treats them as exposure diagnostics rather
than alpha claims. The delta-matched-equities and underlying-Markowitz rows are spanned by
the regressors by construction; their reported \(0.000\) alphas carry numerical-artifact
\(t\)-statistics from dividing by near-zero standard errors and are excluded from any
inference.

\subsection{Robustness reading}

Two robustness cuts summarize the appendix evidence. Leaving out any single underlying
moves the VIX-enabled Greek Markowitz gross Sharpe within a narrow \(1.311\)--\(1.476\)
band around the baseline \(1.374\), while excluding META, NVDA, and TSLA jointly lowers it
to \(1.115\) (Table~\ref{tab:loo_app}); the result is not a one-name artifact, but it does
lean on the volatile mega-cap complex. Across SPY-return terciles, the VIX overlay earns
its keep in down markets (Sharpe \(0.888\) down, \(0.747\) flat, \(2.421\) up, versus
\(-0.973\) down for the equity-option sleeve without the overlay), and the delta-matched
equity benchmark loses in the same states (\(-2.753\) down)
(Table~\ref{tab:regime_performance_app}, Figure~\ref{fig:regime_sharpes_app}). The
walk-forward rolling 36-month diagnostic re-estimates the whole stack each quarter and
still finds a positive out-of-sample Sharpe (Table~\ref{tab:rolling_app}), though on only
twenty evaluation months. The
regime split is descriptive---terciles are formed on the realized sample---but it is the
pattern a volatility-instrument overlay should produce.

\subsection{Claim audit}

\begin{table}[H]
\centering
\scriptsize
\IfFileExists{tables/claim_strength_summary.tex}
  {\resizebox{\textwidth}{!}{\input{tables/claim_strength_summary}}}
  {\begin{tabular}{lll}\toprule Strength & Claim & Boundary \\\midrule Missing & Run the empirical pipeline & -- \\\bottomrule\end{tabular}}
\caption{Claim-strength summary. The table distinguishes durable framework claims from
sample-dependent diagnostics and from live-trading claims that are not made.}
\label{tab:claim_strength}
\end{table}

\begin{table}[H]
\centering
\scriptsize
\IfFileExists{tables/claim_audit.tex}
  {\resizebox{\textwidth}{!}{\input{tables/claim_audit}}}
  {\begin{tabular}{lll}\toprule Claim & Evidence & Status \\\midrule
  Missing & Run the empirical pipeline & blocked \\\bottomrule\end{tabular}}
\caption{Claim-audit summary. The table separates mathematical claims, empirical
diagnostics, settlement-gated VIX claims, post-cost research claims, and production claims
that are not claimed.}
\label{tab:claim_audit}
\end{table}

The claim audit is part of the result. The article can claim a mathematically explicit
option-only efficient frontier and a point-in-time empirical implementation. It can report
gross, exact-settlement VIX, and pre-production post-cost research diagnostics when the
audit layer supports those outputs. It cannot claim production tradability, live fills, live
broker margin parity, or executable capacity because those controls are not broker-run in
this article.
